%% ch01 --- A model is a machine for predicting
%%
%% Written September 2026, as the exemplar every later chapter follows. Every
%% number the reader cannot do in their head comes from code/measure/ch01.py,
%% which imports the SAME rules the listings print (code/ch01/), so the page
%% and the listing cannot disagree about what the model says.
%%
%% WHAT THIS CHAPTER MAY NOT SAY. It plants the book's question -- if the rule
%% is exact and known, is the future known? -- and answers it only for the
%% two gentle models here, where an error in the start dies away or stays the
%% same size. That errors can GROW is Chapter 6's measurement; saying so here
%% as a finding would spend it.
%%
%% Twin: chapters/pl/ch01-models.tex. Section for section, macro for macro,
%% number for number; tools/parity.py compares them.

\chapter{A model is a machine for predicting}
\label{ch:models}
\index{model}

\noindent
You use mathematical models every day without calling them that. The app
that says your train arrives in seven minutes is running one. So is the
oven timer, the tide table, the dosage on a medicine packet and the little
voice that tells you to leave at ten past eight if you want to be there by
nine. Each of them takes a few numbers about now, applies a rule, and hands
you a number about later.

This book is about what happens to that promise. Its first four chapters
show how well it works \dash{} how a line of algebra can tell you where a
stone will be, how hot your tea will be, and when a comet will come back.
The rest of the book shows where it breaks, why it breaks even when the rule
is perfect, and what you can still predict when it does. That last part is
the reason to read it: the breaking is not a failure of mathematics, and
knowing its shape changes what you should believe about every forecast you
are given.

\begin{outcomes}
\outcome{name the state, the rule and the prediction of any model put in
  front of you}
\outcome{run a model written in Python and read what it prints}
\outcome{say what it means for a model to be wrong, and why a wrong model can
  still be useful}
\outcome{say what happens to a small error in the start for the two models in
  this chapter}
\outcome{state the clockwork claim that Laplace made in 1814, and the question
  the rest of the book asks of it}
\end{outcomes}

\section{Three parts of every model}
\label{sec:ch01-parts}
\index{model!state}\index{model!rule}\index{model!prediction}

\noindent
Drop a stone from a bridge. How far has it fallen after three seconds?

A physicist answers with a formula you may remember from school: a stone
falling freely near the Earth's surface covers a distance
\[ d = \tfrac{1}{2}\, g\, t^{2} \]
in $t$ seconds, where $g$ is about \val{ch01.g} metres per second per second.
That one line is a complete model, and it has three parts you will find in
every model in this book.

\begin{itemize}
\item \textbf{The state}: the numbers that say where things are now. Here it
  is a single number, the time $t$ since the stone was let go.
\item \textbf{The rule}: how those numbers turn into the thing you want to
  know. Here it is the formula.
\item \textbf{The prediction}: what the rule says. Here, a distance.
\end{itemize}

\begin{think}
Use the formula for $t = 3$. Roughly how far has the stone fallen: about
$10$ metres, about $30$, or about $45$?
\end{think}
\reveal{About \val{ch01.stone.three} metres: half of \val{ch01.g}, times
$3 \times 3 = 9$.}
If you guessed ten or thirty, you are in good company: most people
underestimate, because the stone does not fall at a steady speed. It keeps
speeding up, so the distance grows with the \emph{square} of the time. Three
times as long a fall is nine times as far.

That is the first thing a model buys you: it catches the intuition that was
going to mislead you. Nobody has a feeling for $t^{2}$. The formula does not
need one.

\begin{trapbox}
\textbf{\enquote{A model is a copy of reality.}} It is not, and nothing about
the formula above resembles a stone. A model is a machine for answering
questions: you put numbers in, it gives numbers out, and the only fair test
of it is whether those numbers agree with what you then measure. A map is
not the territory, and a model is not even a map \dash{} it is a map that
also tells you the time you will arrive.
\end{trapbox}

\section{Your laboratory}
\label{sec:ch01-lab}
\index{Python}\index{laboratory}

\noindent
Every number in this book that you could not work out in your head was
computed by a small program, and every one of those programs is in the
book's repository under \code{code/}. You do not need to be a programmer to
read them. They are short, they are written to be read, and each is
explained the first time it appears. Here is the stone.

\pyregion{code/ch01/falling_stone.py}{model}{The falling stone as a model: one constant, one rule}{lst:ch01-stone}

\noindent
Read it the way you read the formula. The line that starts \code{G =} sets a
constant. The lines that start \code{def distance} define the rule: give it a
time \code{t} and it gives back \code{0.5 * G * t * t}, which is the formula
with the multiplications written out. The text in triple quotes is a note to
the reader, and everything after a \code{\#} is a comment. The program then
asks the rule the question for every whole second from $0$ to $5$:

\pyregion{code/ch01/falling_stone.py}{table}{Asking the model six questions}{lst:ch01-table}

\noindent
\code{for t in range(6)} means \enquote{do the next line with \code{t} set to
$0$, then $1$, and so on up to $5$}. Run it from the \code{code/} directory
with \code{uv run python ch01/falling\_stone.py} and it prints:

\transcript{ch01-falling-stone}

\noindent
Look at the right-hand column. The stone falls about five metres in the
first second and \val{ch01.stone.five} metres in five seconds. If you have
ever dropped something from a height and counted, that second number feels
too large, and it is not.

\begin{note}
The front matter says how to install the laboratory: one tool, \pkg{uv},
which installs the exact version of Python the book was checked on
(\pypatch) and every package at the exact version printed on the copyright
page. After that, \code{uv run python} followed by a file name runs any
listing in the book. If you would rather not run anything, the output of
every listing is printed beside it; running them is how you find out that
the book is not asking you to take its word for anything.
\end{note}

\begin{lab}{l01_01_stone}{The stone, both ways round}
A model answers questions in more than one direction. Given the time, it
gives the distance; given the distance, it gives the time; and given a
measured drop, it tells you what gravity must be. The starter has the three
functions with their descriptions and no bodies. Fill them in until the test
passes. The second one needs a square root, which Python spells
\code{math.sqrt}.
\end{lab}

\section{A model that steps}
\label{sec:ch01-steps}
\index{model!stepping}\index{Newton's law of cooling}

\noindent
The stone's formula lets you jump straight to any time you like. Most
models do not. Pour a cup of tea at \val{ch01.tea.start}\,°C in a room at
\val{ch01.tea.room}\,°C. Isaac Newton described in 1701 how it cools, and in
the form this book uses his rule says: \emph{each minute, the tea loses a
fixed fraction of the gap between itself and the room}. Take the fraction to
be \val{ch01.tea.k}, a tenth.

\begin{think}
The gap at the start is $90 - 20 = 70$ degrees. What is the temperature one
minute later?
\end{think}
\reveal{\val{ch01.tea.one}\,°C. A tenth of the gap is $7$, and $90 - 7 = 83$.}
To know the temperature after two minutes you do the same thing again, to
$83$: the gap is now $63$, a tenth of it is \num{6.3}, and the tea is at
\num{76.7}. There is no shortcut you are expected to know. The rule tells you
how to get from one minute to the next, and to get anywhere else you take one
step at a time.

\pyregion{code/ch01/cooling_cup.py}{rule}{The rule of the next minute}{lst:ch01-rule}

\noindent
The function \code{next\_minute} is the whole model. Given the temperature
now, it returns the temperature a minute later. To see the future you call it
again and again, feeding each answer back in:

\pyregion{code/ch01/cooling_cup.py}{run}{Feeding each answer back in, thirty times}{lst:ch01-run}

\transcript{ch01-cooling-cup}

\noindent
The line \code{temp = next\_minute(temp)} is the most important line in this
book, and you will see it in every chapter under a different name. It
replaces the state with the next state. Doing that over and over is called
\emph{iteration}, and Chapter~\ref{ch:iteration} is about nothing else.

\begin{think}
From the printout, the tea is above $60$\,°C at minute $5$. Is it below
$60$ at minute $6$? What would you have to do to be sure?
\end{think}
\reveal{Yes: it first drops below $60$ at minute \val{ch01.tea.drinkable}. To
be sure you take one more step with the rule, or ask the program for every
minute rather than every fifth.}
That is how you ask a stepping model a question: you step until the answer
appears. It is also why stepping models were rare before computers and are
everywhere now. A machine does not mind taking a million steps.

\plotfig{ch01-two-models}{Two models of the same kind of question. The stone's
formula gives any time directly; the tea's rule gives the next minute, and the
curve is its steps joined up.}{a formula and a stepping rule}

\begin{lab}{l01_02_cooling}{How long until I can drink it?}
Write \code{cool}, which runs the rule for a number of minutes and keeps every
temperature, and \code{minutes\_until}, which answers the question a person
actually asks. Your answer for a cup at $90$\,°C, a room at $20$ and a
fraction of \num{0.1} should agree with the one above.
\end{lab}

\section{Wrong, and useful}
\label{sec:ch01-wrong}
\index{model!error}

\noindent
Neither model is true. The stone's formula ignores the air, and a feather
dropped beside the stone shows you how much that can matter. The tea's rule
ignores the cup, the saucer, the steam and your spoon. Measure carefully and
you will find both are wrong.

\begin{think}
Since the stone's formula ignores the air, and air obviously matters, is the
formula useless?
\end{think}
\reveal{No. For a stone falling a few seconds it is right to within the width
of a hand; for a feather it is badly wrong. A model is useful for the
questions it answers well enough, and it is your job to know which those
are.}

\begin{trapbox}
\textbf{\enquote{A model that is wrong is useless.}} Every model is wrong,
because every model leaves something out \dash{} that is what makes it small
enough to use. The statistician George Box's often-quoted line says the same
thing more briefly: all models are wrong, but some are useful. The useful
question is never \enquote{is it true?} but \enquote{is it good enough, for
this question, over this range?} Keep that question. Part~II of this book is
about a range over which a model with no error at all in its rule stops being
good enough, and the reason is not what you would guess.
\end{trapbox}

\noindent
You can check a model against a measurement, and the check can run in either
direction. Drop a stone from a ten-metre wall and the formula says it lands
after \val{ch01.stone.ten.time} seconds. Time a real one and get a slightly
different number, and you can turn the formula round and ask what $g$ would
have to be. That is the third function in the first lab above, and it is how a model
earns trust: by making predictions sharp enough to be wrong, and then not
being.

\section{What happens to a small mistake?}
\label{sec:ch01-error}
\index{error!in the start}

\noindent
Here is the question the whole book turns on, asked of the gentlest model in
it. You never know the start exactly. Your thermometer says $90$\,°C; the tea
is really at $91$. How wrong is your prediction half an hour later?

\begin{think}
Two cups start one degree apart and cool by the same rule. After thirty
minutes, is the gap between them bigger than one degree, still one degree, or
smaller?
\end{think}
\reveal{Smaller: about \val{ch01.gap.thirty}\,°C. Each minute the gap is
multiplied by \val{ch01.gap.factor}, because each cup loses a tenth of its
own distance from the room.}

\pyregion{code/ch01/two_cups.py}{gap}{Two cups, one degree apart}{lst:ch01-gap}

\transcript{ch01-two-cups}

\noindent
The error does not merely stay small; it disappears. Every minute it is
multiplied by \val{ch01.gap.factor}, so after ten minutes it is about a
third of a degree and after \val{ch01.gap.hundredth} minutes it is less than
a hundredth of one. On a logarithmic scale, where each step down is a factor
of ten, the gap falls along a straight line:

\plotfig{ch01-gap-shrinks}{The gap between two cups that started one degree
apart. On a logarithmic scale it falls along a straight line: the same
fraction is lost every minute.}{an error that dies away}

\noindent
Keep that straight line in your head. A model whose errors shrink forgives
you for not knowing the start. The stone's model is nearly as kind: get the
moment of release wrong by a tenth of a second and your prediction is off by
a fixed amount of time for ever, never by more. Most of the mathematics you
met at school describes systems like these, and it is why the idea that
nature is predictable feels like common sense.

\begin{worldbox}
The dosing instructions on a medicine packet rest on a model of this kind.
Most drugs are cleared from the blood by a roughly fixed fraction every hour,
exactly as the tea loses a fixed fraction of its gap (alcohol is the famous
exception), so a dose taken a little late or a little large leaves an error
that shrinks rather than grows. That forgiveness is why a schedule can be
printed on a packet at all.
\end{worldbox}

\section{The clockwork}
\label{sec:ch01-clockwork}
\index{Laplace's demon}\index{Halley's comet}\index{clockwork universe}

\noindent
Models like these did something astonishing in the eighteenth century.
Edmond Halley, using Newton's laws of motion and gravity, argued that the
comets seen in 1531, 1607 and 1682 were one object on one orbit, and
predicted that it would return around 1758. He died in 1742. The comet was
seen again on Christmas night 1758, and it has carried his name ever since.

A rule on paper had reached seventy-six years into the future and been
right. Pierre-Simon Laplace drew the conclusion in 1814, in a passage that
reads, in its classic English translation:

\begin{quote}
Given for one instant an intelligence which could comprehend all the forces
by which nature is animated and the respective situation of the beings who
compose it \dash{} an intelligence sufficiently vast to submit these data to
analysis \dash{} it would embrace in the same formula the movements of the
greatest bodies of the universe and those of the lightest atom; for it,
nothing would be uncertain and the future, as the past, would be present to
its eyes.
\end{quote}

\noindent
The imagined intelligence is known as \emph{Laplace's demon}, and the picture
behind it as the \emph{clockwork universe}: know the rule and the present,
and the future follows, as surely as the hands of a clock. In 1960 the
physicist Eugene Wigner wrote of \enquote{the unreasonable effectiveness of
mathematics in the natural sciences}, and he meant the same wonder: rules
written by people at desks keep turning out to be right about the world.

\mermaidfig{ch01-model-loop}{Every model in this book has this shape. A state
goes in, a rule turns it into a prediction, and a measurement decides whether
to trust it.}{state, rule, prediction, check}

\begin{think}
Suppose the rule is exact \dash{} not approximately right, but perfectly
right \dash{} and you know it. Can you then predict the future as far ahead as
you like?
\end{think}
\reveal{Only if you also know the start exactly, and you never do. For the
tea and the stone that hardly matters, because the error in the start
shrinks or stays the same size. Whether it always behaves that kindly is the
question the rest of this book answers.}

\noindent
Most people answer yes, and so did most scientists for two hundred years
after Laplace. Hold on to your answer. In Chapter~\ref{ch:logistic} you will
meet a rule one line long, simpler than the tea's, whose errors do not shrink
at all; in Chapter~\ref{ch:butterfly} you will measure how fast they grow; and
by the end of the book you will be able to say, with numbers you computed,
why some things can be predicted for a century and others not for a
fortnight. That is chaos theory, and it is not a theory of disorder. It is
the mathematics of exactly this question.

\begin{summarybox}
\begin{itemize}
\item Every model has three parts: a \textbf{state} (the numbers that say
  where things are now), a \textbf{rule} (how they change, or how they turn
  into an answer), and a \textbf{prediction} (what the rule says).
\item Some models are formulas you can evaluate at any time you like. Most
  are \textbf{stepping rules}: they tell you how to get from now to the next
  moment, and you reach the future by applying them over and over. That is
  iteration.
\item Every model is wrong, because every model leaves something out. The
  useful question is whether it is good enough for this question over this
  range.
\item For the cooling tea, an error in the start \textbf{shrinks} by the same
  factor every minute; for the stone it stays the same size. Models like these
  forgive you for not knowing the start exactly.
\item Laplace's demon is the clockwork claim: know the rule and the present
  exactly, and the future follows. The rest of the book tests it.
\end{itemize}
\end{summarybox}

\canyou

\begin{checkpoint}
\item Name the state, the rule and the prediction of a bus timetable.
  \answerto{The state is the time now and the stop you are at; the rule is
  the timetable (together with the assumption that the bus keeps to it); the
  prediction is when the next bus arrives.}
\item How far does a stone fall in the first two seconds, air ignored? And in
  four?
  \answerto{About \num{19.6} metres in two seconds and about \num{78.5} in
  four. Twice the time, four times the distance.}
\item Tea at $80$\,°C cools in a room at $20$\,°C, losing a tenth of the gap
  each minute. What is its temperature after one minute, and after two?
  \answerto{$74$\,°C and then \num{68.6}\,°C: the gap goes from $60$ to $54$
  to \num{48.6}.}
\item Two cups start two degrees apart in the same room and cool by that rule.
  How far apart are they after one minute?
  \answerto{\num{1.8} degrees: each cup loses a tenth of its own gap to the
  room, so their difference loses a tenth too.}
\item Why can a model that is known to be wrong still be useful?
  \answerto{Because every model leaves something out, and usefulness is a
  question of whether its answers are close enough for the question being
  asked, over the range it is used for.}
\item What would Laplace's demon need to know to predict the future, and which
  of those does the rest of this book say can never be had exactly?
  \answerto{The rule and the exact present state. The rule may be known
  exactly; the present state never is, because every measurement has an
  error.}
\end{checkpoint}
